\documentclass{article}
\usepackage[utf8]{inputenc}
\usepackage[a4paper, left=1.25in, right = 1.25in]{geometry}
\usepackage{graphicx}
\usepackage{wrapfig}
\usepackage{setspace}
\usepackage{titling}
\usepackage{csquotes}
\usepackage{comment}
\usepackage{multicol}
\usepackage{hyperref}

\let\oldsection\section
\renewcommand\section{\clearpage\oldsection}

\setlength{\droptitle}{-5em}   

\def\changemargin#1#2{\list{}{\rightmargin#2\leftmargin#1}\item[]}
\let\endchangemargin=\endlist

\begin{document}
% =-=-= Cover Page Start =-=-=
% \input{Upper Division/Front Page}
% =-=-= Cover Page End =-=-=

% =-=-= Question 1 Start =-=-=
\section{Pony Express}

After getting his first internship at USPS, Renegade came into office from lunch break one day to find out that the postal carriers decided to have a water balloon fight in every server room breaking each rack. In a scramble to keep their jobs, the devs worked to rewrite the back-end. Renegade was told to ensure the validity of 22-digit tracking numbers. Unfortunately, he sprained his foot in his fountain toss and needs your help implementing the validation.  \\
\\
The indices of the digits in the code are 1-based, with 22 being the left-most digit and 1 being the right-most. The check-sum digit is in position 1. The check digit must follow a general rule to be valid. The check digit is the smallest number which, when added to sum of the odd-digits (not-including the check-digit) times three plus the sum of even-digits, gives a number that is a multiple of 10. Note that an ID only contains positive numbers.



\subsection{Input}

The input consists of an arbitrary number of lines, each containing a 22-Digit tracking number.

\subsection{Output}

Each line of the output should contain the validity of a tracking number in the form of true/false.

\subsection{Sample Input/Output}

% 6780726047469215599210
% 2302318695420634188792

\begin{center}
	\begin{tabular}{ |l|l| }
		\hline
		\textbf{Sample Input} & \textbf{Sample Output} \\
	   	\hline
		6780726047469215599210 & true\\
		2302318695420634188792 & false\\
        ... & ...\\
        4638852541045574473209 & true\\
        \hline
	\end{tabular}
\end{center}
% =-=-= Question 1 End =-=-=

\end{document}
